% Generated draft: requires mathematical and source review.
\documentclass[11pt]{article}
\usepackage[margin=25mm]{geometry}
\usepackage{fontspec}
\setmainfont{Latin Modern Roman}
\usepackage{amsmath,amssymb,mathtools,unicode-math}
\setmathfont{Latin Modern Math}
\usepackage{enumitem,needspace,xurl,hyperref}
\setlength{\parindent}{0pt}
\setlength{\parskip}{5pt}
\setlength{\emergencystretch}{3em}
\newcommand{\N}{\mathbb{N}}
\newcommand{\Z}{\mathbb{Z}}
\newcommand{\Q}{\mathbb{Q}}
\newcommand{\R}{\mathbb{R}}
\newcommand{\C}{\mathbb{C}}
\newcommand{\F}{\mathbb{F}}
\newcommand{\A}{\mathbb{A}}
\newcommand{\PP}{\mathbb{P}}
\newcommand{\E}{\mathbb{E}}
\newcommand{\eps}{\varepsilon}
\newcommand{\dd}{\,\mathrm{d}}
\begin{document}
\section*{UnsolvedMath ID: 30001707}
\textbf{Source problem code:} catalogue entry
\subsection{Definitions and mathematical statement}
The supplied record concerns coadjoint orbits and quantization in a noncompact setting, but does not specify the admissible groups, subgroups, or quantization procedure. The following records its questions with these qualifications unresolved.

For the usual interpretation of the notation, let \(G\) be the underlying Lie group with Lie algebra \(\mathfrak g\), and let \(H\) be a Lie subgroup with Lie algebra \(\mathfrak h\). A coadjoint orbit \(O_G\subseteq\mathfrak g^*\) is an orbit for the dual of the adjoint action of \(G\); similarly, \(O_H\subseteq\mathfrak h^*\) denotes an \(H\)-coadjoint orbit. The usual candidate for the map in the record is
\[
\mu:O_G\longrightarrow\mathfrak h^*,\qquad
\mu(\xi)=\xi|_{\mathfrak h}.
\]
The record does not explicitly define \(\mu\), so this interpretation requires confirmation. With this interpretation, define
\[
n(O_G,O_H):=\#\bigl(\mu^{-1}(O_H)/H\bigr),
\]
where the quotient denotes the set of \(H\)-orbits in \(\mu^{-1}(O_H)\), and \(\#\) denotes cardinality. Thus \(n(O_G,O_H)\leq1\) means that this preimage is empty or consists of a single \(H\)-orbit; no finiteness assumption is imposed in defining the count.

Write \(Q(O_G)\) for the representation assigned to \(O_G\) by the intended quantization procedure, and \(Q(O_G)|_H\) for its restriction to \(H\). Where irreducible multiplicities are defined, multiplicity-free means that every such multiplicity is at most one, with the appropriate almost-everywhere interpretation if a direct-integral decomposition is intended. The record does not specify which representation category or decomposition convention applies.

For the intended class of data, the first question is whether
\[
Q(O_G)|_H\text{ is multiplicity-free}
\quad\Longleftrightarrow\quad
\bigl(\forall\,O_H\subseteq\mathfrak h^*\text{ an }H\text{-coadjoint orbit}\bigr)\;
n(O_G,O_H)\leq1.
\]
The universal quantifier on the right ranges over every \(H\)-coadjoint orbit, as stated in the supplied record, without a supplied restriction to quantizable or integral orbits.

For a family \((O_G^\lambda)_{\lambda\in\Lambda}\), write \(\mu_\lambda\) for the corresponding map and
\[
n(O_G^\lambda,O_H)
:=\#\bigl(\mu_\lambda^{-1}(O_H)/H\bigr).
\]
The second question, reading multiplicity-freeness of the restrictions as a simultaneous hypothesis on the family, is whether
\[
\begin{gathered}
\bigl(\forall\lambda\in\Lambda\bigr)\;
Q(O_G^\lambda)|_H\text{ is multiplicity-free}
\\[2pt]
\Longrightarrow
\\[2pt]
\bigl(\forall\lambda\in\Lambda\bigr)
\bigl(\forall\,O_H\subseteq\mathfrak h^*\text{ an }H\text{-coadjoint orbit}\bigr)\;
n(O_G^\lambda,O_H)\leq1.
\end{gathered}
\]
Neither \(\Lambda\), the construction of the family, nor the precise scope of the multiplicity-freeness hypothesis is supplied. Consequently, these are qualified formulations of the source questions, not an assertion that either statement holds for arbitrary Lie groups and quantizations.
\subsection{Short English statement}
The questions ask whether multiplicity-free restriction of a quantized coadjoint orbit is equivalent to every relevant inverse image containing at most one subgroup orbit, and whether multiplicity-free restrictions throughout a family force the same bound at every parameter. The source quantifies over all subgroup coadjoint orbits. Its group hypotheses, quantization, decomposition convention, and family parameters remain unspecified.
\subsection{Sources}
\begin{itemize}
\item Supplied local UnsolvedMath registry record, permanent ID \(30001707\), title \emph{Multiplicity-Free Quantization and Orbit Counts}, source problem code \texttt{catalogue entry}, in the Oberwolfach Reports Open Problems collection. The record supplies the proposed year as 2011 and names no individual proposer.
\item Supplied report citation: \emph{Geometric Quantization in the Non-compact Setting} (2011), Oberwolfach Reports, pp.\ 425--521, DOI \texttt{10.4171/owr/2011/09}. Source URL and supplied literature-source link, labeled ``original Oberwolfach report'': \url{https://doi.org/10.4171/owr/2011/09}. The supplied background describes Oberwolfach Reports as documenting mathematical workshops and their open-problem discussions. No specific problem page, contribution, or individual attribution is supplied.
\item Literature assessment dated exactly \texttt{2026-08-21}: the supplied curation labels the problem open and states that the report presents special symmetric-pair cases and evidence, while the stated noncompact multiplicity-free conjecture is not proved in general. This is supplied dated curation, not independently verified; the record itself requests independent verification of time-sensitive status and cited claims.
\item The registry labels the statement \texttt{corrected\_verified} and reports manual semantic, typographic, and LaTeX review of a meaning-preserving cleanup. These are supplied editorial assertions, not independent verification of the report, mathematical hypotheses, or current problem status.
\end{itemize}
\end{document}
